\loeskopf{Einheit 1 von 5 · Wurzelziehen und Lösbarkeit (Nr. 1--9)}

\erg{1}{a) linear \quad b) quadratisch \quad c) quadratisch \quad d) linear \quad
e) quadratisch}

\erg{2}{a) zwei \quad b) keine \quad c) eine \quad d) zwei \quad e) keine}

\erg{3}{a) $x_1 = 3$, $x_2 = -3$ \quad b) $x_1 = 8$, $x_2 = -8$ \quad
c) $x_1 = 11$, $x_2 = -11$ \quad d) $x_1 = 15$, $x_2 = -15$ \quad
e) $x_1 = \sqrt{11} \approx 3{,}32$, $x_2 = -\sqrt{11} \approx -3{,}32$ \quad
f) keine Lösung \quad g) $x^2 = 100$; $x_1 = 10$, $x_2 = -10$ \quad
h) $x^2 = 16$; $x_1 = 4$, $x_2 = -4$ \quad i) $x^2 = 25$; $x_1 = 5$, $x_2 = -5$ \quad
j) $x - 5 = \pm 6$; $x_1 = 11$, $x_2 = -1$ \quad
k) $x + 3 = \pm 9$; $x_1 = 6$, $x_2 = -12$ \quad l) $x = 7$ (genau eine Lösung)}

\erg{4}{a) wahr ($64 = 64$) \quad b) wahr ($(-4)^2 = 16$) \quad
c) falsch ($5^2 = 25 \neq 16$)}

\erg{5}{a) zwei Lösungen, $x_1 = 2$, $x_2 = -2$ \quad b) eine Lösung, $x = 0$ \quad
c) keine Lösung (die Gerade $y = -5$ liegt unter dem Scheitel)}

\erg{6}{Fehler: das Vorzeichen der $6$ wurde nicht umgedreht. Aus der Klammer folgt
$x + 6 = 5$ oder $x + 6 = -5$, also $x = 5 - 6$ bzw. $x = -5 - 6$. Richtig:
$x_1 = -1$, $x_2 = -11$. \quad a) $x + 4 = \pm 9$; $x_1 = 5$, $x_2 = -13$}

\erg{7}{a) Ein Quadrat ist nie negativ -- keine Zahl mal sich selbst ergibt $-20$.
\quad b) Nein. Nur die Null ergibt beim Quadrieren null, also gibt es genau eine
Lösung, $x = 0$. „Keine Lösung`` und „Lösung null`` sind nicht dasselbe. \quad
c) Die Zahl rechts ist positiv; dann gibt es immer zwei Lösungen. Sie sind nur keine
ganzen Zahlen, sondern $\sqrt{7}$ und $-\sqrt{7}$.}

\erg{8}{a) $x^2 = 64$, $x = 8$; nur die positive Lösung ist eine Länge. Seitenlänge
$8\,$m. \quad b) $x^2 = 200$, $x = \sqrt{200} \approx 14{,}14$; $14{,}14\,$m
$> 14\,$m, das Grundstück passt nicht. \quad
c) $r^2 = \frac{V}{\pi \cdot h} = \frac{50}{\pi \cdot 2{,}5} \approx 6{,}37$;
$r \approx 2{,}52\,$m}

\erg{9}{a) wahr (die Klammer wird nur für $x = 5$ null) \quad b) falsch: aus
$x + 9 = \pm 4$ folgt $x_1 = -5$ und $x_2 = -13$ \quad c) zum Beispiel $x^2 = -9$
oder $x^2 + 4 = 0$ (jede Gleichung, bei der $x^2$ eine negative Zahl sein müsste)}
